\subsection{All-degree principal angles of the original theta quotient sections}\label{endpoint-all_degree_angle_review--all-degree-principal-angles-of-the-original-theta-quotient-sections}

13 September 2026. This is a complete independent derivation of the nonlinear arithmetic determinant comparison in every quotient degree. All spaces are the original polynomial spaces, all projections use the actual common interpolated source Gram, and every change of representative is exhibited as a theta boundary. The earlier sealed AT and AW sources are unchanged.

\subsubsection{1. Fixed original objects and the canonical projection relation}\label{endpoint-all_degree_angle_review--fixed-original-objects-and-the-canonical-projection-relation}

Retain a full actual finite zero packet ENDPOINTSOURCED3545BE978ACCA48978D9119SLOT0END for ENDPOINTSOURCED3545BE978ACCA48978D9119SLOT1END, an integer ENDPOINTSOURCED3545BE978ACCA48978D9119SLOT2END, its original tensor-sum annihilator ENDPOINTSOURCED3545BE978ACCA48978D9119SLOT3END of degree ENDPOINTSOURCED3545BE978ACCA48978D9119SLOT4END, and

ENDPOINTSOURCED3545BE978ACCA48978D9119SLOT5END

Use the original increasing-power remainder frame of ENDPOINTSOURCED3545BE978ACCA48978D9119SLOT6END and the original source coordinate ENDPOINTSOURCED3545BE978ACCA48978D9119SLOT7END. The two source densities are exactly

ENDPOINTSOURCED3545BE978ACCA48978D9119SLOT8END

Their literal masses ENDPOINTSOURCED3545BE978ACCA48978D9119SLOT9END and ENDPOINTSOURCED3545BE978ACCA48978D9119SLOT10END remain in all Grams. Put ENDPOINTSOURCED3545BE978ACCA48978D9119SLOT11END, ENDPOINTSOURCED3545BE978ACCA48978D9119SLOT12END, and let ENDPOINTSOURCED3545BE978ACCA48978D9119SLOT13END be its common Gram on ENDPOINTSOURCED3545BE978ACCA48978D9119SLOT14END. The source proofs establish positivity almost everywhere and every polynomial moment for the two densities. Therefore ENDPOINTSOURCED3545BE978ACCA48978D9119SLOT15END and every restricted Gram below are positive definite. Inner products are conjugate-linear in their first argument, with ENDPOINTSOURCED3545BE978ACCA48978D9119SLOT16END.

For ENDPOINTSOURCED3545BE978ACCA48978D9119SLOT17END, retain the exact sequence

ENDPOINTSOURCED3545BE978ACCA48978D9119SLOT18END

Let ENDPOINTSOURCED3545BE978ACCA48978D9119SLOT19END be the common-metric orthogonal projection onto ENDPOINTSOURCED3545BE978ACCA48978D9119SLOT20END, and ENDPOINTSOURCED3545BE978ACCA48978D9119SLOT21END the projection onto ENDPOINTSOURCED3545BE978ACCA48978D9119SLOT22END. In particular ENDPOINTSOURCED3545BE978ACCA48978D9119SLOT23END. Since ENDPOINTSOURCED3545BE978ACCA48978D9119SLOT24END, the two projections commute and ENDPOINTSOURCED3545BE978ACCA48978D9119SLOT25END. Hence

ENDPOINTSOURCED3545BE978ACCA48978D9119SLOT26END

Polynomial division gives ENDPOINTSOURCED3545BE978ACCA48978D9119SLOT27END. The original minimum section ENDPOINTSOURCED3545BE978ACCA48978D9119SLOT28END satisfies

ENDPOINTSOURCED3545BE978ACCA48978D9119SLOT29END

These objects retain their dependence on ENDPOINTSOURCED3545BE978ACCA48978D9119SLOT30END. We suppress it only inside pointwise equations.

For any two cutoffs ENDPOINTSOURCED3545BE978ACCA48978D9119SLOT31END, both sections have remainder equal to their input. Thus ENDPOINTSOURCED3545BE978ACCA48978D9119SLOT32END has image in ENDPOINTSOURCED3545BE978ACCA48978D9119SLOT33END. Meanwhile ENDPOINTSOURCED3545BE978ACCA48978D9119SLOT34END lies in ENDPOINTSOURCED3545BE978ACCA48978D9119SLOT35END, orthogonal to that image. Orthogonal projection of ENDPOINTSOURCED3545BE978ACCA48978D9119SLOT36END onto ENDPOINTSOURCED3545BE978ACCA48978D9119SLOT37END consequently gives the exact relation

ENDPOINTSOURCED3545BE978ACCA48978D9119SLOT38END

No equality between the different canonical sections is assumed. The same formula is valid at ENDPOINTSOURCED3545BE978ACCA48978D9119SLOT39END, with zero difference.

Let ENDPOINTSOURCED3545BE978ACCA48978D9119SLOT40END, with every map understood through its original coefficient inclusion in ENDPOINTSOURCED3545BE978ACCA48978D9119SLOT41END. The relation coefficient and its full energy are

ENDPOINTSOURCED3545BE978ACCA48978D9119SLOT42END

Indeed ENDPOINTSOURCED3545BE978ACCA48978D9119SLOT43END, and the splitting in ENDPOINTSOURCED3545BE978ACCA48978D9119SLOT44END shows that this is ENDPOINTSOURCED3545BE978ACCA48978D9119SLOT45END. If the relation domain is zero, the formula uses its unique zero map and no inverse on a nonzero space. This is the actual source relation energy that will become the squared sine data.

\subsubsection{2. Finite principal-angle calculation with all Gram factors}\label{endpoint-all_degree_angle_review--finite-principal-angle-calculation-with-all-gram-factors}

All square roots in this section are the positive Hermitian square roots in the fixed coefficient frame. Define the isometric frames

ENDPOINTSOURCED3545BE978ACCA48978D9119SLOT46END

They satisfy ENDPOINTSOURCED3545BE978ACCA48978D9119SLOT47END; this equality retains the original Grams rather than assigning new source norms. Their exact cross matrix is

ENDPOINTSOURCED3545BE978ACCA48978D9119SLOT48END

The order of its factors matters. Multiplication by its actual adjoint gives

ENDPOINTSOURCED3545BE978ACCA48978D9119SLOT49END

Its positivity follows from ENDPOINTSOURCED3545BE978ACCA48978D9119SLOT50END, and its upper bound from AA6. Also ENDPOINTSOURCED3545BE978ACCA48978D9119SLOT51END, retaining the complete source boundary energy.

Let ENDPOINTSOURCED3545BE978ACCA48978D9119SLOT52END be an orthonormal eigenbasis of ENDPOINTSOURCED3545BE978ACCA48978D9119SLOT53END, with eigenvalues ENDPOINTSOURCED3545BE978ACCA48978D9119SLOT54END in ENDPOINTSOURCED3545BE978ACCA48978D9119SLOT55END, including every multiplicity. Define

ENDPOINTSOURCED3545BE978ACCA48978D9119SLOT56END

There is no division by zero because of AA10. Since ENDPOINTSOURCED3545BE978ACCA48978D9119SLOT57END, direct calculation gives

ENDPOINTSOURCED3545BE978ACCA48978D9119SLOT58END

For example the second equality follows from ENDPOINTSOURCED3545BE978ACCA48978D9119SLOT59END, divided by ENDPOINTSOURCED3545BE978ACCA48978D9119SLOT60END. The ENDPOINTSOURCED3545BE978ACCA48978D9119SLOT61END form a basis of ENDPOINTSOURCED3545BE978ACCA48978D9119SLOT62END because they are ENDPOINTSOURCED3545BE978ACCA48978D9119SLOT63END orthonormal vectors in that space.

Define ENDPOINTSOURCED3545BE978ACCA48978D9119SLOT64END by ENDPOINTSOURCED3545BE978ACCA48978D9119SLOT65END. These are the complete principal angles of the two actual section spaces, as witnessed by the paired frames AA11--AA12. Their squared cosines are exactly the eigenvalues of AA10, with no interchange of noncommuting square roots.

If ENDPOINTSOURCED3545BE978ACCA48978D9119SLOT66END, then the orthogonal projection ENDPOINTSOURCED3545BE978ACCA48978D9119SLOT67END has the full norm of ENDPOINTSOURCED3545BE978ACCA48978D9119SLOT68END, so ENDPOINTSOURCED3545BE978ACCA48978D9119SLOT69END and ENDPOINTSOURCED3545BE978ACCA48978D9119SLOT70END. Conversely every ENDPOINTSOURCED3545BE978ACCA48978D9119SLOT71END is fixed by ENDPOINTSOURCED3545BE978ACCA48978D9119SLOT72END on ENDPOINTSOURCED3545BE978ACCA48978D9119SLOT73END, and its coefficient vector is an eigenvector of AA10 of eigenvalue one. Hence

ENDPOINTSOURCED3545BE978ACCA48978D9119SLOT74END

including all repeated eigenvalues of one.

For every ENDPOINTSOURCED3545BE978ACCA48978D9119SLOT75END with ENDPOINTSOURCED3545BE978ACCA48978D9119SLOT76END, put ENDPOINTSOURCED3545BE978ACCA48978D9119SLOT77END. Equations AA12 imply that the vectors ENDPOINTSOURCED3545BE978ACCA48978D9119SLOT78END are orthonormal, orthogonal to every ENDPOINTSOURCED3545BE978ACCA48978D9119SLOT79END, and that different two-dimensional spaces ENDPOINTSOURCED3545BE978ACCA48978D9119SLOT80END are orthogonal. They are also orthogonal to every common vector with ENDPOINTSOURCED3545BE978ACCA48978D9119SLOT81END. On each such plane, the difference of the original projections is exactly

ENDPOINTSOURCED3545BE978ACCA48978D9119SLOT82END

This follows because ENDPOINTSOURCED3545BE978ACCA48978D9119SLOT83END is projection onto ENDPOINTSOURCED3545BE978ACCA48978D9119SLOT84END on that plane, while ENDPOINTSOURCED3545BE978ACCA48978D9119SLOT85END is projection onto ENDPOINTSOURCED3545BE978ACCA48978D9119SLOT86END. The trace is zero and the determinant is ENDPOINTSOURCED3545BE978ACCA48978D9119SLOT87END, giving eigenvalues ENDPOINTSOURCED3545BE978ACCA48978D9119SLOT88END. On the common vectors and on ENDPOINTSOURCED3545BE978ACCA48978D9119SLOT89END, the difference is zero. Thus if ENDPOINTSOURCED3545BE978ACCA48978D9119SLOT90END, its full spectrum on ENDPOINTSOURCED3545BE978ACCA48978D9119SLOT91END is

ENDPOINTSOURCED3545BE978ACCA48978D9119SLOT92END

This proves the complete finite angle-to-projection spectral relation, including zero angles, repeated positive angles, and the remaining zero eigenvalues.

Finally, taking determinants of AA10 gives

ENDPOINTSOURCED3545BE978ACCA48978D9119SLOT93END

Every determinant is computed in the unchanged remainder frame. In particular this identity retains the full original quotient metric, not merely its eigenvalue locations.

\subsubsection{3. The two original cutoff pairs and their forced common direction}\label{endpoint-all_degree_angle_review--the-two-original-cutoff-pairs-and-their-forced-common-direction}

Use the original grouped pairs

ENDPOINTSOURCED3545BE978ACCA48978D9119SLOT94END

The second pair is ordered for ENDPOINTSOURCED3545BE978ACCA48978D9119SLOT95END; for ENDPOINTSOURCED3545BE978ACCA48978D9119SLOT96END its two indices both equal one. Apply AA11--AA16 to both pairs, retaining all ENDPOINTSOURCED3545BE978ACCA48978D9119SLOT97END angles, and define the original endpoint quantity

ENDPOINTSOURCED3545BE978ACCA48978D9119SLOT98END

At least one angle in the second pair is exactly zero. To prove this from the original spaces, both ENDPOINTSOURCED3545BE978ACCA48978D9119SLOT99END and ENDPOINTSOURCED3545BE978ACCA48978D9119SLOT100END lie in

ENDPOINTSOURCED3545BE978ACCA48978D9119SLOT101END

They lie in that hyperplane because each is orthogonal to its own relation space, which contains ENDPOINTSOURCED3545BE978ACCA48978D9119SLOT102END. The latter is a nonzero line inside ENDPOINTSOURCED3545BE978ACCA48978D9119SLOT103END, so it cuts that ENDPOINTSOURCED3545BE978ACCA48978D9119SLOT104END-dimensional space by exactly one dimension in a positive form. Since each ENDPOINTSOURCED3545BE978ACCA48978D9119SLOT105END has dimension ENDPOINTSOURCED3545BE978ACCA48978D9119SLOT106END, ENDPOINTSOURCED3545BE978ACCA48978D9119SLOT107END. Now AA13 proves the claim. At ENDPOINTSOURCED3545BE978ACCA48978D9119SLOT108END, the two one-dimensional spaces are identical, in agreement with their identical cutoffs.

Keep the complete angle list. At each ENDPOINTSOURCED3545BE978ACCA48978D9119SLOT109END, remove one of these forced entries ENDPOINTSOURCED3545BE978ACCA48978D9119SLOT110END only for the purpose of writing the following finite sum, and put

ENDPOINTSOURCED3545BE978ACCA48978D9119SLOT111END

The remaining ENDPOINTSOURCED3545BE978ACCA48978D9119SLOT112END entries include any further zero angles and still sum to ENDPOINTSOURCED3545BE978ACCA48978D9119SLOT113END. No continuous choice of the removed entry or of individual eigenvectors is needed: the next inequality is a pointwise symmetric inequality in this entire list. The smooth quantity differentiated below is the original log determinant, not individual angle branches.

\subsubsection{4. Direct trace bound and finite concavity calculation}\label{endpoint-all_degree_angle_review--direct-trace-bound-and-finite-concavity-calculation}

Define the unchanged relative operators

ENDPOINTSOURCED3545BE978ACCA48978D9119SLOT114END

The last identity follows by factoring ENDPOINTSOURCED3545BE978ACCA48978D9119SLOT115END in its actual order. The operator ENDPOINTSOURCED3545BE978ACCA48978D9119SLOT116END is positive self-adjoint for ENDPOINTSOURCED3545BE978ACCA48978D9119SLOT117END, and ENDPOINTSOURCED3545BE978ACCA48978D9119SLOT118END is self-adjoint for ENDPOINTSOURCED3545BE978ACCA48978D9119SLOT119END because ENDPOINTSOURCED3545BE978ACCA48978D9119SLOT120END is Hermitian. Let its extreme eigenvalues be ENDPOINTSOURCED3545BE978ACCA48978D9119SLOT121END, and set ENDPOINTSOURCED3545BE978ACCA48978D9119SLOT122END.

For any unit vector ENDPOINTSOURCED3545BE978ACCA48978D9119SLOT123END, expansion in an orthonormal eigenbasis of ENDPOINTSOURCED3545BE978ACCA48978D9119SLOT124END gives ENDPOINTSOURCED3545BE978ACCA48978D9119SLOT125END, between those two extremes. Apply this fact to the positive and negative eigenvectors in AA15. Each pair contributes at most ENDPOINTSOURCED3545BE978ACCA48978D9119SLOT126END in absolute value. Summing and using the triangle inequality proves directly

ENDPOINTSOURCED3545BE978ACCA48978D9119SLOT127END

No external trace extremum theorem, or commutation with the section projections, is used.

The original determinant variation is ENDPOINTSOURCED3545BE978ACCA48978D9119SLOT128END. Its four original signs give exactly

ENDPOINTSOURCED3545BE978ACCA48978D9119SLOT129END

It follows either by differentiating the finite determinant ratio from the exact sequence, or from the retained AT12. Apply AA22 to the two pairs in AA17. Using AA18 and AA20 yields the full angle-profile bound

ENDPOINTSOURCED3545BE978ACCA48978D9119SLOT130END

For ENDPOINTSOURCED3545BE978ACCA48978D9119SLOT131END, direct differentiation of ENDPOINTSOURCED3545BE978ACCA48978D9119SLOT132END gives

ENDPOINTSOURCED3545BE978ACCA48978D9119SLOT133END

Thus its derivative decreases, so its secant slopes decrease, proving concavity on ENDPOINTSOURCED3545BE978ACCA48978D9119SLOT134END. Iterating the two-point concavity inequality proves the finite equal-weight inequality ENDPOINTSOURCED3545BE978ACCA48978D9119SLOT135END for positive ENDPOINTSOURCED3545BE978ACCA48978D9119SLOT136END. For nonnegative ENDPOINTSOURCED3545BE978ACCA48978D9119SLOT137END, apply it to ENDPOINTSOURCED3545BE978ACCA48978D9119SLOT138END and let ENDPOINTSOURCED3545BE978ACCA48978D9119SLOT139END; continuity at zero, with ENDPOINTSOURCED3545BE978ACCA48978D9119SLOT140END, gives the same result. Therefore AA24 implies

ENDPOINTSOURCED3545BE978ACCA48978D9119SLOT141END

At ENDPOINTSOURCED3545BE978ACCA48978D9119SLOT142END the finite concavity step is an equality. Additional zero angles in higher degree are still present in AA24; AA26 is the stated consequence using only their sum and the guaranteed one zero angle.

\subsubsection{5. Strict positivity from the actual vertical-line moment source}\label{endpoint-all_degree_angle_review--strict-positivity-from-the-actual-vertical-line-moment-source}

Every term in AA18 is nonnegative. Suppose that ENDPOINTSOURCED3545BE978ACCA48978D9119SLOT143END at one value of ENDPOINTSOURCED3545BE978ACCA48978D9119SLOT144END. Then every angle of the first pair is zero, so AA13 makes its two ENDPOINTSOURCED3545BE978ACCA48978D9119SLOT145END-dimensional spaces equal:

ENDPOINTSOURCED3545BE978ACCA48978D9119SLOT146END

The first is ENDPOINTSOURCED3545BE978ACCA48978D9119SLOT147END, because its relation domain is zero. It contains the original constant polynomial one. The second is orthogonal to ENDPOINTSOURCED3545BE978ACCA48978D9119SLOT148END. Thus ENDPOINTSOURCED3545BE978ACCA48978D9119SLOT149END.

Write the original polynomial in full as ENDPOINTSOURCED3545BE978ACCA48978D9119SLOT150END, ENDPOINTSOURCED3545BE978ACCA48978D9119SLOT151END, and define its actual reflected-conjugate polynomial

ENDPOINTSOURCED3545BE978ACCA48978D9119SLOT152END

Its degree is ENDPOINTSOURCED3545BE978ACCA48978D9119SLOT153END, with leading coefficient ENDPOINTSOURCED3545BE978ACCA48978D9119SLOT154END; this sign is retained. On the unchanged integration line, ENDPOINTSOURCED3545BE978ACCA48978D9119SLOT155END, and so

ENDPOINTSOURCED3545BE978ACCA48978D9119SLOT156END

The polynomial product belongs to the literal relation subspace ENDPOINTSOURCED3545BE978ACCA48978D9119SLOT157END. Orthogonality of one to that space would give

ENDPOINTSOURCED3545BE978ACCA48978D9119SLOT158END

a contradiction. The strict inequality is the positive original Gram of the nonzero polynomial ENDPOINTSOURCED3545BE978ACCA48978D9119SLOT159END, and all integrals exist by the retained moments. This argument uses the actual source measure and vertical-line multiplication identity. It is not a claim about every abstract positive coefficient metric. We have proved

ENDPOINTSOURCED3545BE978ACCA48978D9119SLOT160END

The finite positive Grams and their inverses are smooth near the compact path, so ENDPOINTSOURCED3545BE978ACCA48978D9119SLOT161END is smooth and attains a strictly positive minimum on it.

\subsubsection{6. Nonlinear endpoint transport in every quotient degree}\label{endpoint-all_degree_angle_review--nonlinear-endpoint-transport-in-every-quotient-degree}

Let the full ordered positive spectrum of the original relative Gram be

ENDPOINTSOURCED3545BE978ACCA48978D9119SLOT162END

From AA21 its eigenvalues are ENDPOINTSOURCED3545BE978ACCA48978D9119SLOT163END. Differentiating with respect to ENDPOINTSOURCED3545BE978ACCA48978D9119SLOT164END gives ENDPOINTSOURCED3545BE978ACCA48978D9119SLOT165END, so the same extreme indices apply for every ENDPOINTSOURCED3545BE978ACCA48978D9119SLOT166END, including repeated values and values equal to one. Since each ENDPOINTSOURCED3545BE978ACCA48978D9119SLOT167END is the derivative of ENDPOINTSOURCED3545BE978ACCA48978D9119SLOT168END,

ENDPOINTSOURCED3545BE978ACCA48978D9119SLOT169END

For the fixed integer ENDPOINTSOURCED3545BE978ACCA48978D9119SLOT170END, define on ENDPOINTSOURCED3545BE978ACCA48978D9119SLOT171END

ENDPOINTSOURCED3545BE978ACCA48978D9119SLOT172END

The derivative follows by differentiating ENDPOINTSOURCED3545BE978ACCA48978D9119SLOT173END, using the nonnegative real inverse branch of ENDPOINTSOURCED3545BE978ACCA48978D9119SLOT174END. The square root divisor is strictly positive on the original path by AA31. Multiply AA26 by AA34 and integrate the absolute derivative. Equation AA33 gives the all-degree nonlinear comparison

ENDPOINTSOURCED3545BE978ACCA48978D9119SLOT175END

The original values ENDPOINTSOURCED3545BE978ACCA48978D9119SLOT176END and ENDPOINTSOURCED3545BE978ACCA48978D9119SLOT177END are kept. No monotonicity of ENDPOINTSOURCED3545BE978ACCA48978D9119SLOT178END, individual angles, or individual eigenvectors was assumed.

An equivalent endpoint interval retains both branches explicitly. Set

ENDPOINTSOURCED3545BE978ACCA48978D9119SLOT179END

The endpoint inverse-cosh value lies in ENDPOINTSOURCED3545BE978ACCA48978D9119SLOT180END. Since ENDPOINTSOURCED3545BE978ACCA48978D9119SLOT181END has derivative ENDPOINTSOURCED3545BE978ACCA48978D9119SLOT182END for ENDPOINTSOURCED3545BE978ACCA48978D9119SLOT183END,

ENDPOINTSOURCED3545BE978ACCA48978D9119SLOT184END

The lower bound may be zero while the actual endpoint is strictly positive. When ENDPOINTSOURCED3545BE978ACCA48978D9119SLOT185END, ENDPOINTSOURCED3545BE978ACCA48978D9119SLOT186END, the second pair has its one zero angle and AA35 is the previously proved one-dimensional formula. For any higher ENDPOINTSOURCED3545BE978ACCA48978D9119SLOT187END, including the actual quartet dimension, the same proof applies. It establishes no uniform-in-ENDPOINTSOURCED3545BE978ACCA48978D9119SLOT188END bound for ENDPOINTSOURCED3545BE978ACCA48978D9119SLOT189END.

The new estimate implies the earlier bound ENDPOINTSOURCED3545BE978ACCA48978D9119SLOT190END: the derivative of the inverse function ENDPOINTSOURCED3545BE978ACCA48978D9119SLOT191END is ENDPOINTSOURCED3545BE978ACCA48978D9119SLOT192END. Integrating that derivative between the two endpoint ENDPOINTSOURCED3545BE978ACCA48978D9119SLOT193END values and using AA35 proves the assertion. The angle-profile inequality AA24 remains a still finer pointwise estimate, retaining every individual original Gram ratio.

\subsubsection{7. The literal tensor theta boundary of the angle energy}\label{endpoint-all_degree_angle_review--the-literal-tensor-theta-boundary-of-the-angle-energy}

This section supplies the source primitive behind AA7. Keep the original one-leg complex ENDPOINTSOURCED3545BE978ACCA48978D9119SLOT194END in degrees zero and one, with ENDPOINTSOURCED3545BE978ACCA48978D9119SLOT195END, and the original functions satisfying

ENDPOINTSOURCED3545BE978ACCA48978D9119SLOT196END

The last identity follows by Mellin transformation: both sides transform to the exact ENDPOINTSOURCED3545BE978ACCA48978D9119SLOT197END, and Mellin uniqueness on these source spaces identifies them. The original source intertwining ENDPOINTSOURCED3545BE978ACCA48978D9119SLOT198END lets all polynomial differential operators below act on the complex. Its inclusion in the original two-leg complex sends ENDPOINTSOURCED3545BE978ACCA48978D9119SLOT199END to ENDPOINTSOURCED3545BE978ACCA48978D9119SLOT200END in degree zero and is identity in degree one, so the same primitive maps into the actual tau-base complex with its retained two-leg signs.

In the coefficient ring ENDPOINTSOURCED3545BE978ACCA48978D9119SLOT201END, the cyclic definition of ENDPOINTSOURCED3545BE978ACCA48978D9119SLOT202END gives

ENDPOINTSOURCED3545BE978ACCA48978D9119SLOT203END

Indeed the tensor algebra ENDPOINTSOURCED3545BE978ACCA48978D9119SLOT204END is this polynomial ring modulo that ideal, and its unit is the cyclic vector used to define ENDPOINTSOURCED3545BE978ACCA48978D9119SLOT205END. Evaluating ENDPOINTSOURCED3545BE978ACCA48978D9119SLOT206END on the sum operator and that unit is zero by its actual annihilator definition.

The ideal expression can be chosen by a fully specified finite division. Start with ENDPOINTSOURCED3545BE978ACCA48978D9119SLOT207END; for ENDPOINTSOURCED3545BE978ACCA48978D9119SLOT208END, divide ENDPOINTSOURCED3545BE978ACCA48978D9119SLOT209END by the literal monic ENDPOINTSOURCED3545BE978ACCA48978D9119SLOT210END, viewing the other variables as coefficients:

ENDPOINTSOURCED3545BE978ACCA48978D9119SLOT211END

Later divisions do not increase the already restricted degrees in earlier variables, because ENDPOINTSOURCED3545BE978ACCA48978D9119SLOT212END has complex constant coefficients in those variables. Thus the final ENDPOINTSOURCED3545BE978ACCA48978D9119SLOT213END lies in the span of the tensor remainder monomials, which are a basis of the quotient algebra. Its image in that algebra is zero by AA39, so ENDPOINTSOURCED3545BE978ACCA48978D9119SLOT214END. Summing the exact division identities proves

ENDPOINTSOURCED3545BE978ACCA48978D9119SLOT215END

Every original coefficient of ENDPOINTSOURCED3545BE978ACCA48978D9119SLOT216END and ENDPOINTSOURCED3545BE978ACCA48978D9119SLOT217END remains in these finite quotient polynomials. This construction does not replace ENDPOINTSOURCED3545BE978ACCA48978D9119SLOT218END by a product of simple factors or discard any nilpotent multiplicity.

For a polynomial ENDPOINTSOURCED3545BE978ACCA48978D9119SLOT219END in the original sum variable, define the degree-ENDPOINTSOURCED3545BE978ACCA48978D9119SLOT220END tensor cochain

ENDPOINTSOURCED3545BE978ACCA48978D9119SLOT221END

where the ENDPOINTSOURCED3545BE978ACCA48978D9119SLOT222END-th position is ENDPOINTSOURCED3545BE978ACCA48978D9119SLOT223END in the degree-zero space ENDPOINTSOURCED3545BE978ACCA48978D9119SLOT224END, and every other position is the displayed ENDPOINTSOURCED3545BE978ACCA48978D9119SLOT225END in degree one. There are ENDPOINTSOURCED3545BE978ACCA48978D9119SLOT226END degree-one factors before that position. The tensor differential therefore acts on it with sign ENDPOINTSOURCED3545BE978ACCA48978D9119SLOT227END; this cancels the explicit coefficient sign in AA42. It acts trivially on the other degree-one factors. By AA38, by commutation of the polynomial operators with the differential, and then by AA41, one obtains exactly

ENDPOINTSOURCED3545BE978ACCA48978D9119SLOT228END

This proves the primitive with its full tensor signs. In particular, for ENDPOINTSOURCED3545BE978ACCA48978D9119SLOT229END, regard ENDPOINTSOURCED3545BE978ACCA48978D9119SLOT230END from AA7 as its literal polynomial in ENDPOINTSOURCED3545BE978ACCA48978D9119SLOT231END. Then

ENDPOINTSOURCED3545BE978ACCA48978D9119SLOT232END

Thus the angle energy is the Gram energy of this exact original relation, while its cohomology class is zero. The unchanged arithmetic observation still satisfies ENDPOINTSOURCED3545BE978ACCA48978D9119SLOT233END, with ENDPOINTSOURCED3545BE978ACCA48978D9119SLOT234END and the full unit ENDPOINTSOURCED3545BE978ACCA48978D9119SLOT235END. Its further cohomology image is ENDPOINTSOURCED3545BE978ACCA48978D9119SLOT236END in the original ENDPOINTSOURCED3545BE978ACCA48978D9119SLOT237END. No source relation is dropped merely because those observations kill it.

At every nonempty original support label ENDPOINTSOURCED3545BE978ACCA48978D9119SLOT238END, lift the displayed coefficient maps by ENDPOINTSOURCED3545BE978ACCA48978D9119SLOT239END, and send external absence ENDPOINTSOURCED3545BE978ACCA48978D9119SLOT240END to ENDPOINTSOURCED3545BE978ACCA48978D9119SLOT241END. The class of AA44 is the supported zero at the same label; its full source primitive and energy remain as written. This preserves the original cohomology over the tau base throughout the all-degree estimate.

\subsubsection{8. Full review of the root AGT1--AGT19 source}\label{endpoint-all_degree_angle_review--full-review-of-the-root-agt1agt19-source}

The complete corrected \texttt{tau\_all\_degree\_angle\_transport\_20260913.tex} was read at SHA-256 \texttt{cc0c0c008b0532c5f120d1cea76b94e2d08153c87a694ce18775a5189c72d5f8}. It is accepted without a further mathematical correction. The initial source-cutoff sentence now states ENDPOINTSOURCED3545BE978ACCA48978D9119SLOT242END, making its inclusion ENDPOINTSOURCED3545BE978ACCA48978D9119SLOT243END correct on the entire declared range.

The full acceptance covers the following calculations.

\begin{enumerate}
\def\labelenumi{\arabic{enumi}.}
\tightlist
\item
  AGT1--AGT4 retain the actual packet, full Taylor unit, original source and cohomology observations, both full masses, original vertical line, positive finite Grams, and all four signed determinant terms. The bounded source codomain is now exact.
\item
  AGT5--AGT8 agree with the independent derivation AA6--AA16. In the root source the cross matrix is written in the opposite orientation, ENDPOINTSOURCED3545BE978ACCA48978D9119SLOT244END; its adjoint-times-itself is precisely ENDPOINTSOURCED3545BE978ACCA48978D9119SLOT245END. Thus both presentations give the same squared singular values without commuting the square-root factors. The displayed two-dimensional projection block, all unit-cosine common vectors, and the remaining zero spectrum are correct. The primitive behind its relation-valued difference is supplied explicitly by AA38--AA44 above.
\item
  AGT9--AGT11 retain both full angle lists and remove only one exactly proved zero scalar contribution. The common ambient space has dimension ENDPOINTSOURCED3545BE978ACCA48978D9119SLOT246END, including ENDPOINTSOURCED3545BE978ACCA48978D9119SLOT247END. The trace bound is the sum of the paired positive and negative Rayleigh bounds, so it needs no mutual commutation and introduces no extra factor two.
\item
  AGT12--AGT13 use the precise involution with variable ENDPOINTSOURCED3545BE978ACCA48978D9119SLOT248END and conjugated coefficients. Its leading sign ENDPOINTSOURCED3545BE978ACCA48978D9119SLOT249END, its square, and its equality to complex conjugation on the original line are correct. The resulting integral is the full positive norm of the nonzero original polynomial ENDPOINTSOURCED3545BE978ACCA48978D9119SLOT250END, proving strict positivity on the actual moment path.
\item
  AGT14--AGT19 have the exact constant ENDPOINTSOURCED3545BE978ACCA48978D9119SLOT251END. The derivative and tangent-line proof include the remaining zero-angle values by continuity and the actual positive mean ENDPOINTSOURCED3545BE978ACCA48978D9119SLOT252END. The factor ENDPOINTSOURCED3545BE978ACCA48978D9119SLOT253END in the derivative of ENDPOINTSOURCED3545BE978ACCA48978D9119SLOT254END cancels the factor ENDPOINTSOURCED3545BE978ACCA48978D9119SLOT255END in the angle estimate. The integrated width equals the full original ENDPOINTSOURCED3545BE978ACCA48978D9119SLOT256END, including repeated eigenvalues and endpoint eigenvalues equal to one. Both inverse hyperbolic branches and the explicit endpoint interval are correct. The source states the valid all-degree application, including actual quartet quotients, and preserves the absence of a new uniform estimate for ENDPOINTSOURCED3545BE978ACCA48978D9119SLOT257END.
\end{enumerate}

The separate complete PA1--PA29 derivation in \texttt{tau\_principal\_angle\_volume\_independent\_review\_20260913.md}, SHA-256 \texttt{5f6a75afdbdfd4dcd81781a336702ea6d6545fd8fd0a9b2e014a605e414b3198}, was read in full. Its principal bases, projection spectra, positivity argument, and integrated constant agree with AA1--AA37. This provides an additional independent derivation, with its source scope retained; no finite numerical test is being substituted for these proofs.

\subsubsection{9. Reading, source pins, and mathematical scope}\label{endpoint-all_degree_angle_review--reading-source-pins-and-mathematical-scope}

The complete sealed AT source \texttt{tau\_arithmetic\_determinant\_transport\_20260913.tex} at SHA-256 \texttt{24cf7ce5be6a43cc656641b516185df10b2e4d7824600ece657120d1c549deae} and the complete AW source \texttt{tau\_relation\_window\_spectral\_transport\_20260913.tex} at SHA-256 \texttt{c79e76c087efd59871e3234ff44fb18964515dbd34f0318dffba6b3f1440f6ec} were read in the preceding independent calculation. Their original projection and moment conventions are retained here. The full original \texttt{arithmetic\_input.tex} and \texttt{Tau\_Base\_Cohomology/NOTE.md} were read earlier in this same audit and supply the stated theta complex, finite-packet unit, and tau-base tensor realization. AA6--AA44 provide the new finite angle, analytic transport, and explicit primitive proofs.

The companion receipt pins the final sources and the complete root-source acceptance above. These are closed written proofs and do not depend on a forthcoming supplement. No numerical or Lean execution, source normalization, remote mutation, or unproved uniform eigenvalue bound is claimed. The all-degree conclusion applies to the actual original moment Grams and includes every admissible quotient degree ENDPOINTSOURCED3545BE978ACCA48978D9119SLOT258END.
